As AI systems become more powerful and are deployed more autonomously, it will be increasingly important to prevent them from causing harm. Research in this direction has provided us with a wide array of techniques attempting to mitigate these outcomes, for instance via using models to review outputs of other models, or red-teaming techniques to surface failure modes of current systems.
\\
However, in high-stakes scenarios, it is usually reasonable to prepare and index towards worst-case events. This approach, often called ``black swan'' planning or ``high-impact, low-probability'' event preparation, is crucial in various domains, such as disaster preparedness and national security. In these fields, speculative thinking about potential future threats has proven to be a successful strategy. It is thus only fitting to evaluate whether safety techniques we have in place for AI systems still hold if the systems are themselves intentionally trying to subvert them, regardless of whether we ascribe intentionality or not to them. This is the idea behind AI Control protocols \citep{greenblatt2023ai}. An AI control protocol is a plan for usefully deploying AI systems that prevents an AI from intentionally causing unacceptable outcomes, even if it was actively trying to do so.
\\
Previous work in this sphere evaluates protocols by attempting subversion using an AI following a human-written strategy. Both the safety protocols and the subversion strategy are human provided and only AI executed, effectively limiting the extent of the AI’s influence in the unrolling of the evaluation. Assuming that models will execute best-known strategies will often be overly conservative in the short run and excessively permissive in the long one. In the near future, models will not be capable of optimal responses to control measures, allowing us to learn this from capability evaluations, while in the more distant future, it is not guaranteed that human-specified protocols are going to continue being optimal.
\\
With this in mind, we can leverage adversarial frameworks to jointly train subversion and oversight policies, simulating protocols and strategies in vitro and designing control evaluations that are progressively accounting for model capabilities. We do so by making use of a variant of Checkability Training from \citet{kirchner2024prover}, which closely mimics the Prover-Verifier Game framework of \citet{anil2021learning}. This allows us to model a game where two systems (the ``provers'') engage adversarially to complete a task specification (a ``query''), while being evaluated by a third system (the ``verifier''). The two provers have opposing goals: on one hand, the ``honest prover'' is tasked to complete the task specification correctly. On the other, the ``sneaky prover'' is tasked to complete the task specification introducing an error. Proposed solutions are then provided to the verifier, which outputs a convincingness score. In the helpful role, we reward the prover if it is convincing and correct. In the sneaky role, we reward it for being convincing and incorrect. This setup, borrowed from complexity theory, when designed carefully, gives us as a byproduct a robust monitor for Control evaluations and a subtle, non-human designed subversion policy.
